\cotefile{83}
\scanpage{3}
\[
\alpha^2\deleted{+}+\underbrace{\alpha+1}_{\alpha'}=0
\]

% slanted
\margin{\unsure{Icos-aèdre|Icosaèdre} sur $\mathbb{F}_4$ \uline{gribouillis}}

\[
\begin{aligned}
u^2&=u+\unsure{2|\alpha}\\
u^4&=u^2+\alpha^2=u^{2}+\alpha'=u+\alpha+\alpha'=u+1=u'
\end{aligned}
\]

$u^5=uu'=\alpha\,1$

$u$ d'ordre $15$

\[
\deleted{\unsure{u^3}}\qquad u/u'
\]

\[
u^3=u^2+\alpha u=\underline{u+\alpha}+\alpha u=\alpha'u+\alpha
\]

$u^3$ d'ordre cinq

$u^5$ d'ordre $3$

$\deleted{\unsure{u^3=u'{}^3}}$

$\boxed{\alpha'u+\alpha}$ \deleted{$\unsure{\alpha'u+\alpha}$} opère librement sur $S\smallsetminus\{u,u'\}$ avec $2$ orbites

On veut \qquad $u$

\[
\det(u-v)=0\qquad\text{i.e}\quad \underbrace{\det u}_{1}+\underbrace{\det v}_{1}+\operatorname{Tr}uv=0
\]

\[
\text{-}\qquad\text{i.e.}\quad \operatorname{Tr}uv=0\qquad x\otimes x
\]

\uline{NB}

\[
\begin{aligned}
&\operatorname{Tr}uu'=0\\
&\operatorname{Tr}u.u=\operatorname{Tr}u^2=\operatorname{Tr}(u+\alpha.1)=\operatorname{Tr}u=1\ne0
\end{aligned}
\]

\[
\det(u-v)=\det u\unsure{-|+}\det v+\operatorname{Tr}uv=\operatorname{Tr}uv
\qquad \operatorname{Tr}uv=0
\]

$\deleted{\unsure{\underbrace{u-v}_{1}\quad\underbrace{u-v}_{1}|\det(u-v)}}$

\[
\operatorname{Tr}uv=0\qquad
\operatorname{Tr}(u\underline{v'})=\operatorname{Tr}(u(v+1))=\operatorname{Tr}uv+\operatorname{Tr}u=1+\operatorname{Tr}uv
\]

\figure{\begin{tikzpicture}\draw[rotate=55] (0,0) ellipse (.55 and .22);\end{tikzpicture}}

$\operatorname{Tr}uv$ et $\operatorname{Tr}uv'$ sont \uline{différents}

donc \quad $\operatorname{Tr}uv=0\Longrightarrow\operatorname{Tr}uv'\ne0$

\margin{$\operatorname{Tr}uu'=0$\\$\operatorname{Tr}u^2=1$}

\mspage{L 22}
% typescript upside down relative to the calculations
Tout ce que je croyais avoir à dire, à un moment ou à un autre\deleted{\unsure{d|t}} de la réflexion, que ce soit sur moi, ou sur tel \deleted{autre}\handwritten{\inserted{de mes collègues, \deleted{\illegible{1}} ou amis}}, ou sur un milieu, ou sur\unsure{,} une époque, je l'ai dit, sans \inserted{jamais} avoir à bousculer mes réticences. \handwritten{\inserted{Pour celles-ci,}} Il a suffi à chaque fois que \deleted{\unsure{je les}}\handwritten{\inserted{je les}} examine avec attention, \deleted{\handwritten{\unsure{elles-mêmes}}} pour qu'elles s'évanouissent\deleted{\unsure{,}} sans plus laisser de traces.

% archive sheet number: 2
